% !TeX program = lualatex
% !TeX encoding = utf8
% !TeX spellcheck = uk_UA
% !BIB program = biber


%% --------------------------------------------------------
\chapter{Від рівнянь Айнштайна до рівнянь Фрідмана}
\label{sec:friedmann}
%% --------------------------------------------------------

Космологія описує Всесвіт як єдиний фізичний об'єкт. Її відправний
пункт~--- принцип космологічної однорідності та ізотропії: на
достатньо великих масштабах ($\gtrsim 100\,\mathrm{Mpc}$) розподіл
матерії та геометрія простору-часу не виділяють жодного напрямку і
жодної точки. Цей принцип фіксує єдиний клас метрик~---
Фрідмана--Леметра--Робертсона--Вокера (FLRW)~--- і зводить
рівняння Айнштайна до двох звичайних диференціальних рівнянь для
одного невідомого: масштабного фактора $a(t)$.

Метрика FLRW у \emph{космологічному часі}~$t$ для плоского Всесвіту
($k = 0$):
\begin{equation}\label{eq:FLRW}
  ds^2 = -dt^2 + a(t)^2\,\delta_{ij}\,dx^i dx^j.
\end{equation}
Підставляючи~\eqref{eq:FLRW} у рівняння Айнштайна
$G_{\mu\nu} + \Lambda g_{\mu\nu} = 8\pi G\,T_{\mu\nu}$
для ідеальної рідини
$T^{\mu\nu} = \mathrm{diag}(\rho,\,p,\,p,\,p)$
і обчислюючи ненульові компоненти тензора Айнштайна
\begin{equation*}
  G^0{}_0 = -3H^2,
  \qquad
  G^i{}_j = -(2\dot H + 3H^2)\,\delta^i{}_j,
  \qquad
  H \equiv \frac{\dot a}{a},
\end{equation*}
отримуємо \textbf{рівняння Фрідмана}:

\begin{keybox}
\textbf{Перше рівняння Фрідмана} ($(00)$-компонента):
\begin{equation}\label{eq:F1}
  H^2 = \left(\frac{\dot a}{a}\right)^2
      = \frac{8\pi G}{3}\,\rho + \frac{\Lambda}{3}
\end{equation}
\textbf{Друге рівняння Фрідмана} ($(ii)$-компонента, «прискорення»):
\begin{equation}\label{eq:F2}
  \frac{\ddot a}{a} = -\frac{4\pi G}{3}\,(\rho + 3p) + \frac{\Lambda}{3}
\end{equation}
\textbf{Рівняння збереження} ($\nabla_\mu T^{\mu\nu} = 0$):
\begin{equation}\label{eq:conservation_low}
  \dot\rho + 3H(\rho + p) = 0
\end{equation}
\end{keybox}

\noindent
\textit{Зауваження:}
\begin{itemize}
  \item Рівняння збереження не є незалежним: воно автоматично
        слідує з двох рівнянь Фрідмана через тотожність Б'янкі
        $\nabla_\mu G^{\mu\nu} = 0$.
  \item Космологічна стала $\Lambda$ еквівалентна ідеальній рідині
        з $\rho_\Lambda = \Lambda/(8\pi G)$, $p_\Lambda = -\rho_\Lambda$,
        тобто $w = -1$.
  \item Відхилення від плоскості ($k = \pm 1$) додає член
        $-k/a^2$ у перше рівняння; Planck обмежує
        $|\Omega_k| < 0.005$, тому $k = 0$~--- відмінне наближення.
\end{itemize}

%% --------------------------------------------------------
\section{Рівняння стану і закони розширення}
%% --------------------------------------------------------

Кожна компонента Всесвіту характеризується \textbf{рівнянням стану}
$w \equiv p/\rho$. Рівняння збереження~\eqref{eq:conservation_low}
при сталому $w$ інтегрується точно:
\[
  \dot\rho + 3\frac{\dot a}{a}(1+w)\rho = 0
  \;\Rightarrow\;
  \rho \propto a^{-3(1+w)}.
\]
Підставляємо у перше рівняння Фрідмана $H^2 \propto \rho$:
\[
  \left(\frac{\dot a}{a}\right)^2 \propto a^{-3(1+w)}
  \;\Rightarrow\;
  a(t) \propto t^{\,\frac{2}{3(1+w)}}.
\]
Для випромінювання $\rho \propto a^{-4}$: додатковий множник $a^{-1}$
порівняно з матерією~--- це червоне зміщення фотонів ($E_\gamma \propto 1/a$).

\begin{center}
\begin{tblr}{
  colspec={lllll},
  row{1} = {c, m, font=\bfseries}
}
\toprule
Епоха & Джерело & $w$ & $\rho(a)$ & $a(t)$ \\
\midrule
Інфляція    & Інфлатон              & $\approx -1$ & $\approx\mathrm{const}$ & $\propto e^{Ht}$ \\
Радіаційна  & Фотони, нейтрино      & $1/3$        & $\propto a^{-4}$        & $\propto t^{1/2}$ \\
Матеріальна & Темна матерія+баріони & $0$          & $\propto a^{-3}$        & $\propto t^{2/3}$ \\
Сьогодні    & Темна енергія         & $\approx -1$ & $\approx\mathrm{const}$ & $\propto e^{Ht}$ \\
\bottomrule
\end{tblr}
\end{center}

%% --------------------------------------------------------
\section{Червоний зсув}
%% --------------------------------------------------------

Спостережна кінематика космосу пов'язана з $a(t)$ через
\textbf{космологічний червоний зсув}~$z$. Фотон, випромінений
у момент $t_e$, до спостерігача ($t_0$) приходить розтягненим:
\begin{equation}\label{eq:redshift}
  1 + z \equiv \frac{\lambda_{\rm obs}}{\lambda_{\rm em}}
             = \frac{a(t_0)}{a(t_e)}.
\end{equation}
Нормуємо $a(t_0) \equiv 1$, тоді $a = 1/(1+z)$.
Характерні значення:
\begin{align*}
  z &= 0\ \text{(сьогодні)}, \\
  z &\approx 0.3\ \text{(прискорення розширення)}, \\
  z &\approx 1100\ \text{(рекомбінація)}, \\
  z &\sim 10^{10}\ \text{(нуклеосинтез)}.
\end{align*}

\textbf{Закон Хаббла.}
Фізична відстань між двома точками в просторі FLRW зростає з часом
через розширення:
\[
  d(t) = a(t)\,r,
\]
де $r$~--- незмінна відстань між точками в координатах,
«вморожених» у розширення.
Диференціюючи по часу:
\begin{equation}\label{eq:hubble_law}
  \dot{d} = \dot{a}\,r = \frac{\dot{a}}{a}\,d = H\,d.
\end{equation}
Тобто дві галактики віддаляються одна від одної зі швидкістю,
пропорційною відстані між ними~--- це і є \textbf{закон Хаббла}.
Сьогодні ($a = 1$): $\dot{d} = H_0\,d$. Сучасне значення:
\[
  H_0 = 100\,h\ \mathrm{км/с/Мпк}, \qquad h \approx 0.674\ \text{(Planck\,2018)}.
\]
При малих $z \ll 1$ розкладаємо $a(t) \approx 1 - H_0(t_0 - t_e)$,
тому $z \approx H_0(t_0 - t_e)$ і $d \approx c(t_0 - t_e)$, звідси:
\[
  v \approx cz \approx H_0\,d.
\]
У космології через розширення немає єдиного поняття «відстань»~---
залежно від того, \emph{що} вимірюється, використовують різні визначення.
Базова~--- відстань вздовж світлового конуса:
\begin{equation}\label{eq:comoving_dist}
  d_C(z) = \frac{c}{H_0}\int_0^z \frac{dz'}{E(z')},
  \qquad E(z) \equiv \frac{H(z)}{H_0},
\end{equation}
де $H(z)$ задається формулою~\eqref{eq:Hz}.
З неї будуються спостережні відстані:
\begin{itemize}
  \item \textbf{Відстань люмінозності} $d_L = d_C(1+z)$.
        У пласкому нерозширюваному просторі яскравість джерела
        $J = L / 4\pi r^2$.
        У розширюваному Всесвіті фотони червоніють ($E \to E/(1+z)$)
        і розтягується часовий інтервал між ними~--- обидва ефекти
        дають по множнику $(1+z)$, тому $J = L / 4\pi d_L^2$.
        Саме $d_L$ визначають зі спостережень наднових типу~Ia.
  \item \textbf{Кутова відстань} $d_A = d_C/(1+z)$.
        Об'єкт фізичного розміру $\ell$, що видно під кутом $\delta\theta$,
        знаходиться на відстані $d_A = \ell/\delta\theta$.
        Саме $d_A$ визначає положення акустичних піків CMB
        (§\,\ref{sec:sachs_wolfe_main}).
\end{itemize}

\begin{keybox}
\textbf{Hubble tension.}
Різниця між двома типами відстаней лежить в основі однієї з найгостріших
сучасних космологічних суперечностей.
Локальні методи (наднові типу~Ia, цефеїди~--- вимірюють $d_L$) дають
$H_0 \approx 73\,\mathrm{км/с/Мпк}$.
Ранні методи (кутовий масштаб піків CMB~--- вимірюють $d_A$) дають
$H_0 \approx 67.4\,\mathrm{км/с/Мпк}$.
Розбіжність перевищує $5\sigma$ і може вказувати на нову фізику
за межами $\Lambda$CDM.
\end{keybox}


\begin{figure}[h!]
\centering
\begin{tikzpicture}[scale=0.75]
  % кутова відстань (адаптовано з cosmology.tex)
  \def\p{2}
  \draw[gray!40] (0,0) circle (\p);
  \draw[accent] (0,0) -- node[right, font=\small] {$d_A$} (75:\p)
                (0,0) -- (105:\p);
  \draw[theoryred, ultra thick] (0,0)++(105:\p) arc (105:75:\p)
    node[above, pos=0.5, font=\small, black] {$\ell$};

  \def\p{4}
  \draw[gray!40] (0,0) circle (\p);
  \draw[gray!60] (0,0) -- (75:\p)
                 (0,0) -- (105:\p);
  \draw[theoryred!40, ultra thick] (0,0)++(105:\p) arc (105:75:\p)
    node[above, pos=0.5, font=\scriptsize, black] {$d_C\,\delta\theta$};

  \def\p{0.7}
  \draw[accent!60] (0,0)++(105:\p) arc (105:75:\p)
    node[above, pos=0.5, font=\small] {$\delta\theta$};

  % --- ТЕКСТ ВЗДОВЖ КОЛА ---

  % Мале коло: момент випромінювання (радіус трохи більший за 2, щоб текст був ззовні)
  \path[decorate, decoration={
    text along path,
    text={|\scriptsize\color{gray}|момент випромінювання},
    text align=center
  }] (-50:2.2) arc (-50:-20:2.2);

  % Велике коло: сьогодні (радіус трохи більший за 4)
  \path[decorate, decoration={
    text along path,
    text={|\scriptsize\color{gray}|сьогодні},
    text align=center
  }] (-50:4.2) arc (-50:-20:4.2);

\end{tikzpicture}
\caption{Кутова відстань $d_A = \ell/\delta\theta$.
  Об'єкт фізичного розміру $\ell$ видно під кутом $\delta\theta$;
  через розширення $d_A = d_C/(1+z) < d_C$.}
\label{fig:angular_distance}
\end{figure}


%% --------------------------------------------------------
\section{Параметри густини і модель $\Lambda$CDM}
%% --------------------------------------------------------

Зручно нормувати густину кожної компоненти на
\textbf{критичну густину}~--- значення, за якого Всесвіт плаский:
\begin{equation}\label{eq:rho_crit}
  \rho_{\rm cr}(t) \equiv \frac{3H^2(t)}{8\pi G},
  \qquad
  \Omega_i \equiv \frac{\rho_i}{\rho_{\rm cr}}.
\end{equation}
Тоді перше рівняння Фрідмана~\eqref{eq:F1} переписується як:
\[
  \sum_i \Omega_i = 1 \quad (k=0),
\]
а параметр Хаббла через $z$:
\begin{equation}\label{eq:Hz}
  H(z) = H_0\sqrt{\Omega_r(1+z)^4 + \Omega_m(1+z)^3 + \Omega_\Lambda},
\end{equation}
де кожен член~--- внесок своєї компоненти з таблиці вище.

Стандартна космологічна модель $\Lambda$CDM описує Всесвіт
чотирма компонентами:

\begin{center}
\begin{tblr}{colspec={llll}, row{1}={font=\bfseries}}
\toprule
Компонента & $\Omega_i$ & $w$ & Частка сьогодні \\
\midrule
Баріонна матерія          & $\Omega_b$       & $0$    & ${\approx}5\%$  \\
Холодна темна матерія     & $\Omega_c$       & $0$    & ${\approx}27\%$ \\
Темна енергія ($\Lambda$) & $\Omega_\Lambda$ & $-1$   & ${\approx}68\%$ \\
Випромінювання            & $\Omega_r$       & $1/3$  & ${\approx}0\%$  \\
\bottomrule
\end{tblr}
\end{center}

З погляду спостережної космології, модель повністю параметризується
мінімальним набором із шести чисел\footnote{Фізичні густини баріонів
$\Omega_b h^2$ та темної матерії $\Omega_c h^2$, оптична глибина реіонізації $\tau$,
амплітуда $A_s$ і спектральний індекс $n_s$ первинних збурень,
та кутовий масштаб звукового горизонту $\theta_* \equiv r_s(\eta_*)/d_A(\eta_*)$~---
кут під яким видно звуковий горизонт на поверхні останнього розсіяння
(детально~--- §\,\ref{sec:sachs_wolfe_main}).}.
Їхні найдостовірніші значення зафіксовані місією Planck~2018
\cite{planck_2018}.

%% --------------------------------------------------------
\section{Конформний час і горизонт Хаббла}
%% --------------------------------------------------------

Для опису динаміки Всесвіту часто набагато зручніше перейти від
фізичного (космічного) часу $t$ до \emph{конформного часу} $\eta$,
який визначається як
\begin{equation*}
  d\eta = \frac{dt}{a(t)}.
\end{equation*}
У цих змінних FLRW-метрика набуває
кон\-форм\-но-пласького вигляду:
\begin{equation*}
  ds^2 = a^2(\eta)\bigl(-d\eta^2 + \delta_{ij}\,dx^i dx^j\bigr).
\end{equation*}
У таких координатах розширення Всесвіту повністю «заховане»
у загальному масштабувальному множнику $a(\eta)$. Можна уявити,
що сама координатна сітка розширюється разом із простором.
Через це «адреси» (супутні координати) галактик залишаються незмінними,
а реальна фізична відстань між ними збільшується виключно тому,
що росте масштабний коефіцієнт $a(\eta)$\footnote{Це справедливо лише
для галактик, що беруть участь у загальному космологічному розширенні
(хаббловському потоці). Галактики з власним рухом відносно
нього~--- наприклад, Андромеда (M31), що наближається до
Чумацького Шляху зі швидкістю ${\sim}110\,\mathrm{km/s}$~---
мають координати, що змінюються з часом. На космологічних
масштабах такі \emph{пекулярні швидкості} малі порівняно
з темпом розширення і зазвичай ними нехтують.}.

\begin{keybox}
Перехід до конформного часу є базовим інструментом лінеаризованої теорії
гравітації та квантової космології. По-перше, він усуває недіагональні
компоненти метрики, істотно спрощуючи калібрувальні умови рівнянь Ейнштейна.
По-друге, у конформних змінних безмасові або конформно-інваріантні поля
підпорядковуються хвильовому рівнянню стандартного \textbf{гармонічного осцилятора}
в просторі Мінковського, де еволюція фонової геометрії враховується виключно
через ефективний часозалежний потенціал. Ця властивість робить конформний
час природним формалізмом для опису квантової генерації та еволюції
первинних збурень (розділи \ref{sec:metric}–\ref{sec:quantum}).
\end{keybox}

Позначимо похідну за $\eta$ штрихом і введемо конформний
параметр Хаббла $\mathcal{H} \equiv a'/a = aH$. Рівняння
Фрідмана у конформних змінних:
\begin{equation*}
  \mathcal{H}^2 = \frac{8\pi G}{3}\,a^2\bar{\rho}
                + \frac{\Lambda}{3}\,a^2,
  \qquad
  \mathcal{H}' = -\frac{4\pi G}{3}\,a^2(\bar{\rho} + 3\bar{p})
               + \frac{\Lambda}{3}\,a^2,
\end{equation*}
де $\bar{\rho}(\eta)$ та $\bar{p}(\eta)$~--- однорідні фонові
густина та тиск. Комбінуючи ці два рівняння, отримуємо чисто
кінематичне співвідношення, що зв'язує темп зміни горизонту
з рівнянням стану:
\begin{equation}\label{eq:calH_prime}
  \mathcal{H}' = \mathcal{H}^2\!\left[1 - \tfrac{3}{2}(1+w)\right],
  \qquad w \equiv \frac{\bar{p}}{\bar{\rho}}.
\end{equation}
Знак правої частини~\eqref{eq:calH_prime} визначає динаміку
горизонту: при $w > -1/3$ (стандартне розширення) фізичний
радіус Хаббла $R_H = c/H$ зростає, а при $w < -1/3$
(інфляція)~--- стискається. Фізично $R_H = c/H = H^{-1}$ (у натуральних одиницях $c=1$)~--- це масштаб,
на якому швидкість розширення Всесвіту дорівнює швидкості
світла; у теорії збурень він відокремлює режими, що
коливаються ($k \gg \mathcal{H}$), від «заморожених»
($k \ll \mathcal{H}$)\footnote{Це поняття динамічного
горизонту не слід плутати з горизонтом частинок, який
визначає абсолютну причинну зв'язаність від початку часів.}.

Розв'язуючи~\eqref{eq:calH_prime} для сталого $w$, отримуємо
закон розширення $a(\eta)$ для кожної космологічної епохи:

\begin{center}
\begin{tblr}{
  colspec={lllll},
  row{1} = {c, m, font=\bfseries}
}
\toprule
Епоха & Джерело & $w$ & $a(\eta)$ & $a(t)$ \\
\midrule
Інфляція    & Інфлатон
            & $\approx -1$
            & $\propto -\dfrac{1}{H\eta}\ \ (\eta < 0)$\footnotemark
            & $\propto e^{Ht}$ \\
Радіаційна  & Фотони, нейтрино
            & $1/3$
            & $\propto \eta$
            & $\propto t^{1/2}$ \\
Матеріальна & {Темна матерія \\ + баріони}
            & $0$
            & $\propto \eta^2$
            & $\propto t^{2/3}$ \\
Сьогодні    & Темна енергія
            & $\approx -1$
            & $\propto -\dfrac{1}{H\eta}\ \ \text{(асимптотика)}$
            & $\propto e^{Ht}$ \\
\bottomrule
\end{tblr}
\end{center}
\footnotetext{Під час інфляції $a(t)=a_0 e^{Ht}$, тому
$\eta(t)=\int_{-\infty}^{t}\frac{dt'}{a_0 e^{Ht'}}
=-\frac{1}{a(t)H}=-\frac{1}{\mathcal{H}}$.
Оскільки $a(t)H>0$ завжди, $\eta<0$ протягом усієї інфляції;
зручно покласти $\eta\to 0^-$ на її кінці, тобто вся епоха~---
$\eta\in(-\infty,0)$.
Конформний горизонт Хаббла $\mathcal{H}^{-1}=-\eta=1/(aH)$
при цьому \emph{спадає}: $a$ росте при $H\approx\mathrm{const}$,
що є кінематичним визначенням інфляції.}

%% --- Місток до §2 ---
$\Lambda$CDM є теорією \emph{еволюції}, а не
\emph{початкових умов}: вона не виводить з перших принципів (\emph{ab initio})
ні надзвичайної просторової пласкості Всесвіту, ні природи
і спектру первинних флуктуацій густини~--- ці параметри
($A_s$, $n_s$, про них йтиметься далі) приймаються як феноменологічні вхідні дані,
що визначаються з спостережень. Їхнє фізичне обґрунтування
забезпечує теорія космічної інфляції~--- предмет наступного розділу.
